\section{结论}
本文提出 \our{}，一个以编排为中心的智能体系统。它通过统一四元组接口（Instruction、Context、Tools、Model）自动创建子智能体，以解决复杂长程智能体任务。通过将子智能体视为可动态创建的单元，编排器能够按需生成针对任务定制的执行器，并为其提供专用工作记忆（指令、上下文）和能力（模型、工具）。

这一抽象具有现实价值：它以任务充分的上下文支持按需专门化，并使不同实现下的子智能体保持即插即用。抽象与执行的解耦使编排器具备可学习性；我们给出了两种优化方式，即监督微调与上下文学习。实验结果表明，当搭配 Gemini-3-Flash 等前沿模型时，\our{} 在三个高难度基准（GAIA、Terminal-Bench 和 SWE-Bench-Verified）上均取得强劲且稳定的提升。它显著优于已有基线框架，例如在所有基准上的 pass@1 平均提高 16.28 个百分点。这些结果验证了以编排为中心的方法在自动处理复杂长程任务时的有效性。
